
\subsubsection{2.1 CCNet and Perplexity-Based Quality Filtering}

Wenzek et al. [2019] introduced CCNet, a pipeline for extracting high-quality monolingual text from CommonCrawl by training language-specific KenLM n-gram models on Wikipedia and scoring web documents by perplexity. The original design is explicitly per-language: each language receives its own KenLM model trained on its Wikipedia, and cutoffs are applied from language-specific tercile boundaries stored in \texttt{cutoff.csv}. This per-language calibration acknowledges that perplexity scales depend on the reference corpus and should not be compared across languages.

However, practitioners adapting CCNet for large-scale multilingual pipelines have often simplified this design by applying a single global percentile threshold to the combined multi-language document pool, discarding the per-language calibration in the original design. Dolma \citep{Soldaini2024Dolma} applies Pythia-based perplexity scores; RedPajama-V2 \citep{Weber2024RedPajama} provides pre-computed ccnet\_perplexity as a quality signal metadata field intended for downstream filtering. In each case, the path from pre-computed perplexity to actual data selection --- including whether filtering is applied globally or per-language --- is left to the practitioner. This paper quantifies the consequences of the most natural choice (global k-th percentile).

\subsubsection{2.2 Multilingual Dataset Quality Audits}

Caswell et al. [2021] provide the most systematic audit of per-language quality differences in multilingual web corpora, evaluating over 100 languages in mC4. They document severe quality disparities in low-resource languages and argue that per-language quality evaluation is necessary --- but they do not measure the association between language membership and retention rate under a global threshold, and their findings concern different filtering signals (not ccnet\_perplexity specifically).

Adelani et al. [2023] audit quality filtering for African languages and find that global filters consistently fail for low-resource languages. Niyomugabo et al. [2025] (BhashaKritika) identify that per-language KenLM scoring is necessary even as a prerequisite quality evaluation step. These works motivate per-language approaches but do not provide quantitative Cramér's V measurements for a specific dataset-signal-threshold combination that practitioners can calibrate correction strategies against.

\subsubsection{2.3 Retention Rate Tuning for Multilingual Filtering}

The work closest to ours is Turki et al. [2026], who study cross-lingual quality classifiers and find that retention rate tuning is necessary when applying multilingual quality filters --- globally calibrated classifiers systematically under-retain some language groups. They adopt percentile-based (relative) thresholds computed per regression head as a best practice. However, their setting differs: they work with neural classifiers trained on the mC4 corpus, not with pre-computed KenLM perplexity signals from RedPajama-V2. Their endorsement of per-language percentile calibration provides indirect prior support for the correction direction our work motivates, but does not supply the V baseline required to evaluate it.

Singh et al. [2024] find that per-language filtering strategies produce qualitatively different diversity-quality tradeoffs than English-centric approaches. Chaudhary et al. [2025] explicitly note that ''absolute thresholds lack general validity unless supported by extensive ablation'' and adopt percentile-based filtering --- another endorsement of per-language calibration without a direct measurement of the association baseline.

\subsubsection{2.4 RedPajama-V2 and Quality Signals}

Weber et al. [2024] describe the RedPajama-V2 dataset and its quality signal pipeline. They acknowledge that ''ML-based quality signals have been reported to lead to biases or underrepresent minorities'' and include ccnet\_perplexity as one of 40+ quality signals computed for 5 languages (en, de, fr, es, it). They do not quantify the language-group retention disparity produced by global thresholding on any individual quality signal. This paper fills that gap for ccnet\_perplexity.

\subsubsection{2.5 Position of This Work}

This paper is the first to quantify the language-group retention disparity (Cramér's V) from global ccnet\_perplexity thresholding on RedPajama-V2, across multiple threshold levels, with Holm-corrected significance. The V = 0.40--0.57 baseline established here enables future correction studies to be properly powered and evaluated. The finding that prior estimates underestimate actual V by 25--40\% is an empirical contribution to the literature on multilingual quality filtering calibration.

